% !PRINTED 139
\weberpage{139}

Diese Funktionen genügen nach (9) den Differentialgleichungen
\weq{11}{\begin{aligned}
\frac{dx}{dv} &= yz,\\
\frac{dy}{dv} &= -zx,\\
\frac{dz}{dv} &= -\varkappa^2xy,
\end{aligned}}
und den Nebenbedingungen:
\weq{12}{\operatorname{sn}0 = 0, \qquad \operatorname{cn}0 = 1, \qquad \operatorname{dn}0 = 1.}

Es bestehen zwischen ihnen die Relationen
\weq{13}{y^2 = 1 - x^2, \qquad z^2 = 1 - \varkappa^2x^2.}

Aus den Fundamentaleigenschaften der $\vartheta$-Funktionen ergeben
sich die ersten Eigenschaften der elliptischen Funktionen:
\[
\operatorname{sn}v = -\operatorname{sn}(-v), \qquad \operatorname{cn}v = \operatorname{cn}(-v), \qquad \operatorname{dn}v = \operatorname{dn}(-v),
\]
d. h. $\operatorname{sn}v$ ist eine ungerade, $\operatorname{cn}v$, $\operatorname{dn}v$ sind gerade Funktionen.

Setzen wir noch
\weq{14}{\pi\vartheta_{00}^2 = 2K, \qquad \pi\vartheta_{00}^2\omega = 2iK',}
und folglich
\weq{15}{\pi\vartheta_{01}^2 = 2\varkappa'K, \qquad \pi\vartheta_{10}^2 = 2\varkappa K, \qquad \omega = \frac{iK'}{K},}
so erhält $v$ die Werte $K$, $iK'$, $K+iK'$, wenn $u = \frac{1}{2}, \frac{\omega}{2}, \frac{1+\omega}{2}$ wird, und es folgt aus (4) mit Rücksicht auf (3) und \S\,21, (10):
\weq{16}{\begin{aligned}
&\operatorname{sn}K = 1, \quad \operatorname{sn}(K+iK') = \frac{1}{\varkappa},\\
&\operatorname{cn}K = 0, \quad \operatorname{cn}(K+iK') = -\frac{i\varkappa'}{\varkappa},\\
&\operatorname{dn}K = \varkappa', \quad \operatorname{dn}(K+iK') = 0,\\
&\operatorname{sn}iK', \quad \operatorname{cn}iK', \quad \operatorname{dn}iK' = \infty.
\end{aligned}}

Nun lassen sich $K$, $K'$ vermittelst (8) durch bestimmte Integrale
ausdrücken, und man erhält, wenn man $v$ von einem Eckpunkte des
Parallelogramms $0, K, K+iK', iK'$ bis zum folgenden längs der
Peripherie verschiebt, also $u$ längs
\[
0, \qquad \frac{1}{2}, \qquad \frac{1+\omega}{2}, \qquad \frac{\omega}{2}.
\]
